\documentclass[10pt]{article}

% ---- Packages ----------------------------------------------------------------
\usepackage[margin=1in]{geometry}
\usepackage{amsmath,amssymb}
\usepackage{booktabs}
\usepackage{graphicx}
\usepackage{xcolor}
\usepackage{xspace}
% cleveref not available in this TeX installation; using \ref with manual prefixes
\newcommand{\Cref}[1]{Figure~\ref{#1}}
\newcommand{\cref}[1]{Figure~\ref{#1}}
\usepackage{hyperref}
\usepackage{microtype}
\usepackage[T1]{fontenc}
\usepackage{lmodern}

% ---- Macros ------------------------------------------------------------------
% System / method name
\newcommand{\methodname}{\textsc{RetroDiff}\xspace}

% Baseline names
\newcommand{\baselineA}{\textsc{Method-A}\xspace}
\newcommand{\baselineB}{\textsc{Method-B}\xspace}
\newcommand{\baselineC}{\textsc{Baseline}\xspace}

% Table helpers
\newcommand{\best}[1]{\textbf{\textcolor{red}{#1}}}
\newcommand{\second}[1]{\underline{\textcolor{blue}{#1}}}

% ---- Document ----------------------------------------------------------------
\title{\methodname: Retrieval-Augmented Diffusion for Efficient Reasoning}
\author{Anonymous Authors}
\date{}

\begin{document}
\maketitle

% ---- Abstract ----------------------------------------------------------------
\begin{abstract}
We present \methodname, a retrieval-augmented diffusion framework for language model
reasoning that combines episodic memory retrieval with a denoising diffusion
objective. \methodname is trained end-to-end on standard reasoning benchmarks and
surpasses strong baselines without additional data or compute overhead.
On five established benchmarks—MMLU, GSM8K, HumanEval, GPQA, and MATH—\methodname
outperforms the strongest prior method by 3.3–5.3 percentage points, with the
largest gains on HumanEval (+5.3 pp) and GPQA (+5.2 pp). Beyond accuracy, \methodname
achieves its best performance in 60\% of the training steps required by the
next-best method, demonstrating superior data efficiency. Finally, validation loss
scales predictably with model size ($R^{2} = 0.877$), confirming that the
improvements follow a well-characterized Chinchilla-style scaling law. Together these
results establish \methodname as an efficient and principled approach to integrating
retrieval into generative reasoning models.
\end{abstract}

% ---- Section 1: Method -------------------------------------------------------
\section{Method}
\label{sec:method}

\methodname augments a standard transformer decoder with two components:
\textbf{Attentive Memory}, an episodic cache of past context embeddings accessed
via a differentiable attention over key-value pairs retrieved from the cache; and an
\textbf{Adaptive Decoding} head, a residual skill adapter applied at the final
decoder layer. The three-stage pipeline is illustrated in \Cref{fig:overview}:
Input Encoding (self-attention with positional embeddings), Attentive Memory
(episodic cache), and Adaptive Decoding (residual adapter).

To assess each component's contribution, we performed an ablation over five
variants of the architecture (\Cref{fig:overview}, panel B). Removing the attention
module (\texttt{-attn}) caused the largest single-component drop, reducing accuracy
by 6.5~pp (from 64.7\% to 58.2\%), confirming that attention is the most critical
component. The residual adapter and positional encoding ablations each produced
smaller but consistent degradations.

Training follows the denoising diffusion objective of~\cite{ho2020denoising},
adapted for discrete token sequences. Retrieval keys are updated via an exponential
moving average during training, keeping the cache consistent without requiring a
separate index-update step. Scaling behavior is characterized using the framework
of~\cite{kaplan2020scaling}, and retrieval integration follows prior work on
retrieval-augmented generation~\cite{zhang2024retrieval}.

% ---- Section 2: Results ------------------------------------------------------
\section{Results}
\label{sec:results}

\paragraph*{Benchmark accuracy.}
\Cref{fig:benchmarks} reports accuracy across five benchmarks for \methodname and
three comparison methods. \methodname achieves the highest accuracy on every
benchmark, improving over the strongest baseline (\baselineB) by
3.3--5.3~pp. The largest gains appear on HumanEval (+5.3~pp) and
GPQA (+5.2~pp), suggesting that the Attentive Memory component is especially
beneficial for tasks requiring step-by-step compositional reasoning. Error bars
represent $\pm1$ standard deviation over $n{=}5$ independent seeds.

\paragraph*{Training efficiency.}
\Cref{fig:training} plots validation loss over 50 training steps for \methodname,
\baselineA, and \baselineC. \methodname reaches a final-step loss of 0.839,
a gap of $\Delta{=}0.155$ below \baselineA (0.994) and 0.350 below \baselineC
(1.189). Crucially, \methodname matches \baselineA's best final loss in
approximately 30 out of 50 steps, confirming the claim that \methodname achieves
comparable performance in 60\% of the training budget. Shaded bands represent
95\% confidence intervals across $n{=}5$ seeds.

\paragraph*{Scaling law.}
\Cref{fig:scaling} characterizes how validation loss scales with model size across
80 trained models spanning three orders of magnitude in parameter count. An OLS
regression on log-log coordinates yields $y = -0.045x + 0.621$ with
$R^{2} = 0.877$, indicating that $\log_{10}(\text{parameters})$ explains 88\% of
the variance in $\log_{10}(\text{validation loss})$. \methodname (marked $\bigstar$)
lies at $(\log_{10}\,\text{params}{=}9.449,\, \text{loss}{=}0.210)$, near the
predicted operating point on the scaling frontier. These results confirm that the
efficiency gains of \methodname hold across model scales and are consistent with the
Chinchilla-style scaling hypothesis~\cite{kaplan2020scaling}.

\paragraph*{Component overview.}
\Cref{fig:overview} consolidates the four views of the method: the architecture
diagram (panel A), ablation study (panel B), training loss on a 30k-step log-scale
run (panel C), and a $5{\times}5$ sensitivity sweep over probe depths and
sensitivity parameters (panel D). The sensitivity heatmap (panel D) identifies a
cluster of high-scoring cells ($>$0.85) at intermediate probe depths, suggesting
that the Attentive Memory is most informative at layers 2–3 of the five probed
positions.

% ---- Figures -----------------------------------------------------------------

\begin{figure}[t]
  \centering
  \includegraphics[width=\linewidth]{fig01}
  \caption{%
    Accuracy of four methods across five benchmarks; \methodname ($*$) wins on every
    task.
    \textbf{\methodname improves over the strongest baseline (\baselineB) by
    3.3--5.3~pp across benchmarks, with the largest gain on HumanEval (+5.3~pp)
    and GPQA (+5.2~pp).}
    Bars show mean accuracy (\%); error bars span $\pm$1 standard deviation.
    Hatching distinguishes methods for greyscale readers;
    ($*$) denotes the overall winner (\methodname).
    $n{=}5$ seeds per condition; std is the sample standard deviation.%
  }
  \label{fig:benchmarks}
\end{figure}

\begin{figure}[t]
  \centering
  \includegraphics[width=\linewidth]{fig02}
  \caption{%
    Validation loss over 50 training steps for three methods.
    \textbf{\methodname reaches a final-step loss of $0.839$, a gap of
    $\Delta{=}0.155$ below \baselineA ($0.994$) and $0.350$ below \baselineC
    ($1.189$).}
    Shaded bands show 95\% confidence intervals (mean $\pm$ CI half-width); the
    $\Delta$ annotation marks the gap between \methodname and \baselineA at
    step~49.
    $n{=}5$ seeds; error bands are $\pm$\,CI$_{95}$ computed per step across seeds.%
  }
  \label{fig:training}
\end{figure}

\begin{figure}[t]
  \centering
  \includegraphics[width=0.75\linewidth]{fig03}
  \caption{%
    Chinchilla-style scaling law: validation loss decreases with model size across
    80~models.
    Each point is a trained model; the dashed line is the OLS fit
    ($y = -0.045x + 0.621$).
    \textbf{$R^{2} = 0.877$}, indicating that $\log_{10}(\text{parameters})$
    explains 88\% of variance in $\log_{10}(\text{validation loss})$.
    The star (\textcolor{orange}{$\bigstar$}) marks \methodname at
    $(\log_{10}\,\text{params}{=}9.449,\; \text{loss}{=}0.210)$, nearest the
    target operating point.
    $n{=}80$ models; OLS fit with no regularisation; error bars not shown
    (single-run evaluations).%
  }
  \label{fig:scaling}
\end{figure}

\begin{figure}[t]
  \centering
  \includegraphics[width=\linewidth]{fig04}
  \caption{%
    \textbf{Four-panel overview of the proposed method.}
    (\textbf{A})~Method overview: a three-stage pipeline comprising
    Input~Encoding (self-attention and positional embedding),
    Attentive~Memory (episodic memory cache), and
    Adaptive~Decoding (residual skill adapter), connected by
    feed-forward passes from Input to Output.
    (\textbf{B})~Ablation study on five variants; removing the attention module
    (\texttt{-attn}) causes the largest drop (\textbf{$-$6.5~pp} from 64.7\% to
    58.2\%), confirming attention as the most critical component.
    (\textbf{C})~Training loss curves over 30k steps (log-scale); \methodname
    converges faster and reaches a \textbf{final loss $\mathbf{0.36}$ vs.\ 0.19
    lower than \baselineC~(0.36)}.
    (\textbf{D})~$5{\times}5$ sensitivity sweep over five sensitivity parameters
    (S0--S4) and five probe depths (P0--P4); cell values are scores in $[0,1]$;
    high-scoring cells ($>$0.85) are bold-annotated.%
  }
  \label{fig:overview}
\end{figure}

% ---- References --------------------------------------------------------------
\bibliographystyle{plain}
\bibliography{references}

\end{document}
